\documentclass[11pt]{article}
\usepackage[margin=1in]{geometry}
\usepackage{amsmath,amssymb}
\usepackage{xcolor}
\usepackage{enumitem}
\usepackage{hyperref}
\hypersetup{colorlinks=true,linkcolor=blue!50!black,urlcolor=blue!50!black}

\newcommand{\rev}[1]{\textcolor{blue!70!black}{#1}}
\newcommand{\cmt}[1]{\par\smallskip\noindent\textit{#1}\par\smallskip}
\newcommand{\resp}{\par\noindent\textbf{Response.}\ }

\title{Response to Reviewers\\\large Manuscript ToG-2026-0064:\\``From Rules to Nash Equilibria: Formally Verified Game-Theoretic\\Analysis of a Competitive Trading Card Game''}
\author{}
\date{}

\begin{document}
\maketitle

We thank the Editor-in-Chief, the Associate Editor, and the three reviewers for
their careful and constructive reading. The reviewers found the work novel,
technically impressive, and a valuable application of formal verification to game
analysis; their main concern was that several conclusions were stated more
strongly than the modeling assumptions justify, together with requests for
greater accessibility, robustness reporting, and clarity about verification
boundaries. We have revised the manuscript thoroughly to address every point.

\paragraph{How to read the revision.} All substantive additions and changes are
typeset in \rev{blue} in the revised manuscript (a macro \verb|\highlightchangesfalse|
produces a clean camera-ready version). Section references below point to the
revised PDF.

\paragraph{Summary of the most important changes.}
\begin{enumerate}[leftmargin=1.4em]
\item \textbf{Correctness fix to the symmetric-equilibrium numbers.} While
  responding to Reviewer~3 we discovered that the symmetric-Nash gap figures
  reported in the previous version (Dragapult ``$416.7$\textperthousand, gap
  $63.3$,'' ``five best responses'') were inherited from an early, superseded
  enumeration of a \emph{different} game model and did not match the Lean
  artifact that the paper actually verifies (\texttt{SymmetricNash.lean}). The
  verified symmetric equilibrium has value exactly $500$, \emph{seven}-deck
  support, and places Dragapult at $459.6$\textperthousand---a strict gap of
  $\mathbf{40.4}$\textperthousand. We recomputed every affected number with exact
  rational arithmetic, cross-checked it against the Lean theorems, and corrected
  the abstract, Table~VII (and its caption), and Section~VII. We are grateful
  that the review process surfaced this; the corrected figures are now fully
  consistent with the machine-checked artifact.
\item \textbf{Dragapult's exclusion is now shown to be equilibrium-independent}
  (Reviewer~3, Major~\#1--2): a complementary-slackness argument establishes that
  a strict best-response gap against \emph{one} optimal strategy excludes a deck
  from the support of \emph{every} equilibrium---symmetric \emph{or} asymmetric.
\item \textbf{New accessibility material} (Reviewer~2): a Lean/formal-verification
  primer (Sec.~III-A), a notation table (Table~I), a verification-pipeline diagram
  (Fig.~1), and uniform captions on all code listings.
\item \textbf{Foregrounded qualifications} (Reviewer~1): ``suboptimal'' is now
  explicitly model-relative in the abstract, Sec.~X, and the conclusion, and the
  paper is framed as a methodological case study.
\item \textbf{Title reframed as a case study} (Reviewer~1, AE): the title now reads
  ``\ldots A Lean~4 Case Study in Game-Theoretic Analysis\ldots'' instead of
  ``Formally Verified\ldots,'' aligning the headline framing with the more measured,
  case-study presentation requested throughout.
\item \textbf{Artifact counts reconciled exactly with the source} (Reviewer~2). We
  re-audited the published artifact and corrected every count to its exact,
  grep-reproducible value: \textbf{87 files}, \textbf{$\sim$31{,}900 lines},
  \textbf{2{,}627 theorems}, of which \textbf{145} use \texttt{native\_decide} and
  \textbf{2{,}482} are kernel-checked (Table~IX rebuilt with the true per-category
  split; previous version's category counts were stale). One displayed
  expected-win-rate value (Ceruledge) was corrected from $44.8\%$ to its exact
  $45.4\%$. These were the only numeric discrepancies we found on a full recheck.
\item \textbf{The equilibrium-independent exclusion is now machine-backed and
  global} (Reviewer~3). We added a Lean theorem
  (\texttt{symmetric\_nash\_dragapult\_strict\_gap}) certifying Dragapult's strict
  sub-value payoff, and a linear program (\texttt{dragapult\_universal\_exclusion.py})
  that maximizes Dragapult's payoff over the \emph{entire} optimal-strategy set
  ($459.6 < 500$), upgrading the earlier size-bounded search to a global
  certificate. We also replaced a stale supplementary results file that reported
  the superseded ($A{=}M,B{=}M^\top$) model with a corrected enumeration
  (\texttt{enumerate\_symmetric\_nash.py}) consistent with the paper.
\end{enumerate}

\hrulefill

\section*{Associate Editor}
\cmt{``some conclusions are overstated relative to the underlying modeling
assumptions \ldots\ the interpretation of equilibrium results and claims regarding
strategy optimality require more careful qualification \ldots\ dependence on
subjective modeling choices, the limited empirical scope, and the need for clearer
discussion of verification boundaries, robustness, and supporting artifacts \ldots\
improved accessibility.''}
\resp We have acted on each item: (i) equilibrium claims are re-qualified and the
key exclusion result is now proven equilibrium-independent (Sec.~VII);
(ii) ``optimality/suboptimality'' is consistently scoped to the single-match
payoff model (abstract, Sec.~X, conclusion); (iii) the dependence on subjective
type assignments is tested and bounded (Sec.~III-C); (iv) empirical scope is
reframed as a case study (Sec.~X, conclusion); (v) the verification boundary is
classified explicitly as an engineering bottleneck rather than a fundamental
limit (Sec.~IX-A); (vi) robustness is reported via explicit perturbation
($\varepsilon$) bounds (Sec.~X); (vii) accessibility is improved with a primer,
notation table, and pipeline diagram. Details follow per reviewer.

\hrulefill

\section*{Reviewer 1}

\cmt{(1) The paper leans heavily on ``formally verified''/``machine-checked'' in a
way that may overstate assurance, given the reliance on \texttt{native\_decide}.}
\resp We added a measured framing throughout, beginning with the \textbf{title},
which now reads ``\ldots A Lean~4 Case Study\ldots'' rather than ``Formally
Verified\ldots.'' The abstract retains the explicit
trust-boundary sentence and now adds the equilibrium-independence qualifier; the
assurance gradient is summarized in Table~IX (kernel- vs.\ compiler-checked, with
exact counts: $2{,}482$ kernel-checked vs.\ $145$ \texttt{native\_decide}), and
Sec.~IX-A states plainly that \texttt{native\_decide} produces no kernel-checked
proof term. We also classify a fully kernel-checked version as a tractable
engineering task (see Reviewer~3, Minor~\#3).

\cmt{(2) The bridge from rules to empirical conclusions depends on subjective type
assignments; this should be tested more carefully.}
\resp Sec.~III-C now separates load-bearing from borderline assignments and adds a
\emph{quantitative} sensitivity test. Only genuine multi-type decks are borderline;
the single assignment the popularity paradox depends on---Dragapult Dusknoir as a
Psychic defender, weak to the unambiguously Dark counters Grimmsnarl, Mega~Absol,
and N's~Zoroark---is \emph{not} borderline. Enumerating every defensible
reassignment of the borderline decks places the aggregate alignment rate in the
range $[70.6\%,\,87.5\%]$ (point estimate $83.3\%$), all far above the $50\%$
random baseline, while the headline result is invariant across all of them. We now
report this range explicitly and label the alignment figure as a descriptive,
sensitivity-bearing statistic rather than a verified claim.

\cmt{(3) The equilibrium narration is too confident; the exact support is fragile,
and conclusions should distinguish the point estimate from what is stable under
perturbation.}
\resp We now make this distinction central. Sec.~VII and the new ``Equilibrium
robustness'' paragraph (Sec.~X) separate (a) the \emph{perturbation-stable}
exclusion of Dragapult, which tolerates an entrywise perturbation of
$\approx 20$\textperthousand, from (b) the \emph{point-estimate-specific} support
composition, whose marginal members flip under $\approx 5$\textperthousand. Our
headline conclusion rests only on (a).

\cmt{(4) ``Suboptimal'' needs cautious handling; the bimatrix/Swiss/consistency/
familiarity caveat should be brought to the foreground, especially the abstract
and conclusion.}
\resp Done. The abstract now states that ``suboptimal'' is relative to a
single-match payoff model omitting Swiss and consistency incentives; the
conclusion repeats this; and Sec.~X (``Strategic objective mismatch'') spells out
that a player may rationally choose an in-model-suboptimal deck for Swiss
consistency, familiarity/misplay-risk, event stability, or card access.

\cmt{(5) The empirical scope is narrow (three-week window, top-14, third-party
pipeline); present the work as a methodological case study.}
\resp We adopt this framing explicitly. The abstract calls the contribution ``a
reproducible case study''; Sec.~X closes by reading the contribution as a
methodological case study for one well-instrumented format and window rather than
a universal metagame claim; data-provenance biases remain itemized in Sec.~V-C.

\hrulefill

\section*{Reviewer 2 (Minor Revision)}

\subsection*{Accessibility and artifact}
\cmt{The manuscript assumes substantial prior knowledge of formal verification;
the LEAN source was not included in the submission package.}
\resp We added a dedicated primer (Sec.~III-A) introducing proof assistants,
Lean~4, decision procedures, and exact arithmetic for the general AI/PCG audience,
and a pipeline diagram (Fig.~1). On the artifact: the complete anonymized Lean
source (\textbf{87 files, ${\sim}31{,}900$ lines, $2{,}627$ theorems}, no
\texttt{sorry}/\texttt{admit}/axiom---confirmed by \texttt{\#print axioms} on the
headline theorems, which depend only on \texttt{propext}/\texttt{Quot.sound} and
the \texttt{native\_decide} axiom \texttt{Lean.ofReduceBool}), the $14\times14$
matrix, and the Python scripts are provided as supplementary material and re-verify
end-to-end via a single \texttt{lake build}. We re-audited every count against the
source so the manuscript and artifact now agree exactly, removed a superseded
results file and the raw-model enumeration scripts it came from, and added two
reproduction scripts (\texttt{enumerate\_symmetric\_nash.py},
\texttt{dragapult\_universal\_exclusion.py}). The strengthened Data Availability
statement maps each numerical claim to a named theorem so any result can be located
and re-run. The archive is attached to this resubmission.

\subsection*{Clarifications}
\cmt{1. (V-A) Rationale for weighting a tie as $T/3$.}
\resp Added (Sec.~V-A): the coefficient mirrors the Pok\'emon Championship Series
match-point system (win $=3$, tie $=1$, loss $=0$), so a tie is worth exactly
one-third of a win in the format's \emph{own} reward structure---encoded in the
artifact as \texttt{scoreWin}=3, \texttt{scoreDraw}=1, \texttt{scoreLoss}=0. This
is more faithful than the common $T/2$ convention. We also note the
symmetric-equilibrium support is invariant to this coefficient (Sec.~X).

\cmt{2. (IX-B, XI-A) What does the high ``fixed cost'' of porting entail?}
\resp Sec.~IX-B now itemizes it: (i) encoding the strategic-layer rule semantics
(type chart, legality, prize/turn structure)---the bulk of the code and the only
part needing real domain modeling; (ii) defining the archetype enumeration and
matrix ingestion; (iii) reusing the game-theory harness unchanged, since it is
parameterized over an abstract \texttt{FiniteGame}. Steps (ii)--(iii) are days of
work given a matrix; step (i) dominates and scales with rule complexity. A fresh
archetype set within the same game is the ${\sim}200$-line marginal case.

\cmt{3. (VIII) How does the Swiss objective specifically influence the analysis?}
\resp Sec.~VIII now explains that our Nash/replicator results optimize a
single-match criterion, whereas Swiss qualification optimizes a \emph{threshold}
functional $P(\text{X--2})$ over eight rounds. Because the threshold is sensitive
to outcome \emph{variance}, not only the mean, it re-weights the single-match
equilibrium toward consistent decks for strong contenders near the cut-line,
which can rationalize a portion of the observed ``suboptimal'' play. We mark a
full Swiss-utility equilibrium as future work.

\cmt{4. (II) Wilson intervals and perturbations make the method a hybrid,
complementary to Monte Carlo, rather than orthogonal.}
\resp We agree and rephrased Sec.~II accordingly: the approach is now characterized
as a \emph{hybrid} that applies formal verification \emph{on top of} statistical
estimation---Wilson intervals for per-cell uncertainty and a $10{,}000$-iteration
Monte~Carlo resampling for equilibrium robustness---with Lean certifying exact
claims over each sampled or point-estimate matrix. Fig.~1 visualizes this
interface.

\subsection*{Suggestions}
\cmt{1. Inconsistent listing captions.}
\resp All code listings now carry captions (Listings~1--8).
\cmt{2. A notation table would help.}
\resp Added (Table~I), covering $s_j$, $w_{i,j}$, $\mathbb{E}[\mathrm{WR}_i]$, $M$,
$S_{ij}$, the game value $v$, replicator $x_i,f_i,\bar f$, the Wilson $\tilde p$,
$P_{\mathrm{Bo3}}$, the Swiss cut-line, and the permil scale.
\cmt{3. A brief intro to LEAN/formal verification.}
\resp Added as Sec.~III-A (see above).
\cmt{4. A pipeline diagram clarifying the Python\,$\leftrightarrow$\,LEAN interface
and compile-time vs.\ runtime.}
\resp Added as Fig.~1: an untrusted Python discovery layer only \emph{proposes}
candidates; Lean \emph{certifies} every claim at compile time over exact
rationals; a separate runtime stage executes demonstration traces.

\hrulefill

\section*{Reviewer 3 (Major Revision)}

\cmt{Major \#1. Uniqueness via symmetric equilibria is a conceptual leap:
asymmetric equilibria $(A,B)$ are interchangeable and equally rational, so
uniqueness is an artifact of restricting to symmetric profiles, not a consequence
of individual rationality.}
\resp We agree with the theory and have rewritten Sec.~VII to make the stronger,
\emph{equilibrium-independent} argument. We now explicitly acknowledge
interchangeability and that symmetry is a modeling choice. We then show that
Dragapult's exclusion does \emph{not} rely on the symmetry restriction: in a
constant-sum game, complementary slackness forces the support of \emph{every}
optimal strategy to lie inside the best-response set of \emph{every} optimal
opponent strategy. Since the symmetric strategy is optimal for both players and
Dragapult earns $40.4$\textperthousand{} \emph{strictly} below the value against
it, Dragapult is a best response to no optimal strategy and therefore appears in
\emph{no} equilibrium---symmetric or asymmetric. This strict gap is now itself
machine-checked (\texttt{symmetric\_nash\_dragapult\_strict\_gap}), and we
corroborate the universal claim with a \emph{global} linear program
(\texttt{dragapult\_universal\_exclusion.py}) that maximizes Dragapult's payoff over
the \emph{entire} optimal-strategy set and returns $459.6<500$---a stronger
certificate than the earlier size-bounded support search. We now state that the
symmetry restriction only selects a single \emph{canonical} profile to report,
whereas the \emph{exclusion} is a strict best-response fact holding across the
entire equilibrium set; we also delimit the scope precisely (the universal claim is
exact for the constant-sum symmetrization, with the approximately-constant-sum raw
game corroborating via its directly verified zero-weight equilibria).

\cmt{Major \#2. Chance timing/evolutionary interpretation: under an asymmetric
equilibrium the population would split into two sub-populations; the manuscript
implicitly shifts from classical to bounded/evolutionary rationality and should
decouple the two and state the epistemic boundary of the single-population
replicator.}
\resp Added a dedicated paragraph, ``Classical versus evolutionary rationality''
(Sec.~VII-A). We distinguish the classical one-shot equilibrium question from the
evolutionary imitation dynamics, note that an asymmetric profile $(A,B)$ would
physically split the population into two behavioral sub-populations, and state
that the single-population replicator is valid for \emph{macro} share dynamics
only under a specific epistemic regime: an approximately homogeneous population,
largely unstructured matchmaking, and open centralized information flow
(net-decking). We emphasize this is an assumption, not a theorem, and that the
Dragapult-exclusion result lives at the static-equilibrium level and does not
depend on the evolutionary reading.

\cmt{Major \#3. Quantal response/noise: the proof rests on a knife-edge,
noise-free matrix; under payoff perturbations asymmetric equilibria evaporate
while the symmetric profile is a robust focal point. Discuss how the symmetric
equilibrium withstands micro-level payoff uncertainty (Sec.~VI-A).}
\resp Added a paragraph to Sec.~VI-A. We make two points. First, the load-bearing
claim is \emph{not} knife-edge: Dragapult's $40.4$\textperthousand{} gap tolerates
$\approx 20$\textperthousand{} of entrywise perturbation, several times the
$\pm 1.8$--$9$\textperthousand{} Wilson sampling error. Second, noise is precisely
why we anchor on the symmetric profile: asymmetric equilibria of a
near-constant-sum game are highly sensitive to matrix drift, whereas the symmetric
(maximin) strategy is the self-correcting focal point that stochastic,
quantal-response play converges toward; the sensitivity analysis ($>96\%$ core
inclusion) is its empirical signature.

\cmt{Minor \#1. The \texttt{GrimssnarlFroslass} identifier typo; state explicitly
that it is a benign string-matching artifact.}
\resp The footnote now states explicitly that the identifier is a purely cosmetic
string-matching artifact used consistently throughout, with no effect on any
definition, theorem, or result, retained verbatim for traceability.

\cmt{Minor \#2. State the exact $\varepsilon$ bound under which the symmetric Nash
support set remains invariant.}
\resp Added (Sec.~X, ``Equilibrium robustness''). A uniform entrywise perturbation
of size $\varepsilon$ moves any payoff-vs-mix and the value by at most
$\varepsilon$, so a deck with best-response gap $g$ cannot enter the support while
$\varepsilon < g/2$. The smallest exclusion gap among non-support decks is
$10.3$\textperthousand, so an excluded deck can enter only once
$\varepsilon \gtrsim 5$\textperthousand; we now also note the opposite direction (a
thin incumbent such as Alakazam, at $2.5\%$ weight, can exit under a comparable
perturbation), so the single $\varepsilon$ guaranteeing the \emph{entire} support
is ${\approx}5$\textperthousand. Dragapult's exclusion, with gap
$40.4$\textperthousand, is by contrast invariant for $\varepsilon \lesssim 20$\textperthousand{}
(a four-fold larger margin). For the tie coefficient specifically the bound is
exact: because each matchup and its mirror share one record, the tie terms cancel
in $M_{ij}-M_{ji}$, so the symmetrization is algebraically independent of the tie
weight; we verified $\max_{ij}|S_{ij}^{T/3}-S_{ij}^{T/2}|<1$\textperthousand{} and
that re-enumeration under $T/2$ reproduces the identical seven-deck support.

\cmt{Minor \#3. State whether a kernel-checked alternative to \texttt{native\_decide}
is fundamentally impossible or merely an engineering bottleneck.}
\resp Sec.~IX-A now classifies it explicitly as an \emph{engineering bottleneck,
not a fundamental impossibility}: the obstruction is the kernel's inefficient
reduction of \texttt{Fin.foldl}, and a kernel-transparent reimplementation (e.g.,
structurally recursive predicates or inductive relations instead of opaque folds)
is possible in principle at the cost of substantial proof-engineering effort and
longer kernel-reduction times.

\hrulefill

\section*{IEEE AI-content policy}
Per the Editor's note, we added a ``Disclosure of Generative AI Use'' statement:
generative AI assistance was used during revision for copy-editing/reorganization
and for cross-checking numerical computations; all formal results are
machine-checked in Lean~4 and were independently verified by the authors, who take
full responsibility for the content.

\bigskip
We believe these revisions resolve the reviewers' concerns while preserving the
artifact's rigor, and we thank the reviewers again for feedback that materially
strengthened both the paper and its correctness.

\end{document}
